% Generated draft: requires mathematical and source review.
\documentclass[11pt]{article}
\usepackage[margin=25mm]{geometry}
\usepackage{fontspec}
\setmainfont{Latin Modern Roman}
\usepackage{amsmath,amssymb,mathtools,unicode-math}
\setmathfont{Latin Modern Math}
\usepackage{enumitem,needspace,xurl,hyperref}
\setlength{\parindent}{0pt}
\setlength{\parskip}{5pt}
\setlength{\emergencystretch}{3em}
\newcommand{\N}{\mathbb{N}}
\newcommand{\Z}{\mathbb{Z}}
\newcommand{\Q}{\mathbb{Q}}
\newcommand{\R}{\mathbb{R}}
\newcommand{\C}{\mathbb{C}}
\newcommand{\F}{\mathbb{F}}
\newcommand{\A}{\mathbb{A}}
\newcommand{\PP}{\mathbb{P}}
\newcommand{\E}{\mathbb{E}}
\newcommand{\eps}{\varepsilon}
\newcommand{\dd}{\,\mathrm{d}}
\begin{document}
\section*{UnsolvedMath ID: 10900034}
\textbf{Source problem code:} catalogue entry
\subsection{Definitions and mathematical statement}
A hyperbolic \(n\)-simplex is a geodesic simplex in hyperbolic \(n\)-space, with the compact/ideal/other vertex conventions to be checked against the source. A finite set \(\mathcal S\) specifies simplex shapes; finite gluings may use arbitrarily many copies of those shapes. In a facewise isometric gluing, the cone angle around a codimension-two face is the sum of its incident dihedral angles. The source requires the gluing to be a closed cone \(n\)-manifold, not merely any cell complex with assigned angles.

Cooper's Problem 5.9 asks whether
\[
 \forall R>0\ \forall n\in\Z,\ n\ge4\ \exists\eps>0\
 \exists\mathcal S\text{ finite}
\]
with the following two properties:
\begin{enumerate}
\item Every closed cone \(n\)-manifold \(N\) obtained by gluing simplices from \(\mathcal S\), with every codimension-two cone angle \(\theta\) satisfying
\[
 2\pi-\eps<\theta<2\pi+\eps,
\]
admits a hyperbolic metric.
\item Every closed hyperbolic \(n\)-manifold \(M\) with \(\operatorname{inj}_x(M)>R\) at every point is obtained by such a gluing from the same \(\mathcal S\).
\end{enumerate}
Both \(\eps\) and \(\mathcal S\) may depend on \(R,n\), but must work uniformly for every \(M,N\) in these clauses. The first clause asks for existence of a hyperbolic metric on \(N\); it does not assert that the original cone metric is smooth hyperbolic. In the second clause ``obtained this way'' could mean realization of the underlying manifold by an allowed cone gluing, or realization of its given hyperbolic metric. That distinction and the precise simplex and cone-manifold conventions require source review; neither is silently fixed.

The source calls such a collection a Thurston's Lego set and reports existence in dimensions at most \(3\). The supplied assessment reports no verified construction in dimensions at least \(4\), with no complete modern primary-literature audit recorded.
\subsection{Short English statement}
For each radius threshold and dimension at least four, find finitely many hyperbolic simplex shapes that both model every sufficiently thick closed hyperbolic manifold and force any allowed near-\(2\pi\)-angle cone gluing to admit a hyperbolic metric. One common shape set and angle tolerance must serve all the manifolds in both requirements.
\subsection{Sources}
\begin{itemize}
\item Local UnsolvedMath registry, permanent ID \(10900034\), exact source code \texttt{catalogue entry}, subsection 5.9, attributed to Cooper; no proposed year supplied.
\item Delp, Hoffoss and Manning, \emph{Problems In Groups, Geometry, and Three-Manifolds} (2015), 5.9, PDF page 4: \url{https://arxiv.org/pdf/1512.04620#page=4}; supplied status-source URL \url{https://arxiv.org/abs/1512.04620}.
\item The source calls these Thurston's Lego sets and reports their existence in dimensions at most three. No separate lower-dimensional construction reference is supplied.
\item Source-TeX extraction checked against rendered PDF text, accessed 2026-07-29. Source evidence NEEDS\_REVIEW: open in 2015, no complete modern primary-literature audit recorded. Registry status open and classification OPEN-TRIAGE; summary reports no verified primary-literature construction for \(n\ge4\).
\item Dataset attribution: Ulam AI / UnsolvedMath Contributors, CC BY 4.0. Source rights NEEDS\_REVIEW: publicly accessible author-submitted source TeX and PDF, redistribution terms requiring release review.
\end{itemize}
\end{document}
